% title: Least Norwegian Number
% short-title: Least Norwegian Number
% abstract: We formalize in Lean 4 the olympiad-style divisor problem asking for the least positive integer that admits three distinct positive divisors summing to $2022$. The argument is structural: the largest chosen divisor must be the number itself below the critical threshold, and the remaining factorization cases force $1344$.
% subjclass: 11A05
% keywords: divisors, olympiad problem, minimality, Lean 4
% author: Mario Hernández M.

\subsection{Least Norwegian Number}

We define a positive integer to be \emph{Norwegian} if it has three distinct
positive divisors whose sum is $2022$. The formal statement proved in Lean is
\lean{Biblioteca.Demonstrations.no_small_norwegian}, together with
\lean{Biblioteca.Demonstrations.least_norwegian_number}. The first result rules
out every candidate below \(1344\), and the second packages the final minimality
statement.

\begin{theorem}
There is no Norwegian number \(n\) with \(n<1344\).
\end{theorem}

\begin{theorem}
The least Norwegian number is \(1344\).
\end{theorem}

\begin{proof}
The number \(1344\) is Norwegian because \(6\), \(672\), and \(1344\) are
distinct positive divisors of \(1344\), and
\[
  6+672+1344 = 2022.
\]

For the minimality claim, let \(a<b<c\) be positive divisors of a Norwegian
number \(n\) with
\[
  a+b+c = 2022.
\]
Since \(a\le c-2\) and \(b\le c-1\), we obtain \(2022 \le 3c-3\), hence
\(c\ge 675\). If \(c<n\), then \(c\) is a proper divisor of \(n\), so
\(n\ge 2c \ge 1350\), which is impossible when \(n<1344\). Therefore for every
Norwegian number below \(1344\), the largest selected divisor is actually
\(c=n\).

Now write the second divisor as \(b \mid n\), say \(n = db\) with \(d>1\).
Because
\[
  2022 = a+b+n = a + (d+1)b
\]
and \(a<b\), one obtains \(d \le 3\). Thus only the cases \(d=2\) and \(d=3\)
remain. Writing also \(n = am\), the equation \(a+b+n=2022\) becomes
\[
  a(2+3m)=4044 \qquad \text{if } d=2,
\]
and
\[
  a(3+4m)=6066 \qquad \text{if } d=3.
\]
In the first case, the factor \(2+3m\) is congruent to \(2 \bmod 3\), at least
\(11\), and divides \(4044 = 3 \cdot 4 \cdot 337\). The Lean proof uses the
prime factor \(337\) to force \(2+3m = 674\), hence \(m=224\) and \(n=1344\),
contradicting \(n<1344\). In the second case, the same argument with
\(6066 = 18 \cdot 337\) forces \(3+4m=1011\), so \(m=252\) and \(n=1512\),
again impossible.

Thus no Norwegian number lies below \(1344\). Combined with the explicit
witness \(1344\), this proves that \(1344\) is the least Norwegian number. The
corresponding Lean functions are listed in the glossary at the end of the
paper.
\end{proof}
